%% Section 7, Evaluation. Converted from paper/sections/07_evaluation.md.
\section{Evaluation, and a protocol for comparison}
\label{sec:evaluation}

% No reproduced plate in this section: the photograph of the caged motion-capture
% rig that used to stand here illustrated a cost the paragraph below states in
% rollouts and hours, which is the form the reader can use.
% ---------------------------------------------------------------------------
\makeatletter
\@ifundefined{fnLoaded}{\IfFileExists{figs/finding_numbers.tex}{% generated by tools/make_finding_figures.py; do not edit
% Every count the finding figures and their captions print, recomputed from corpus/rows.
\providecommand{\fnLoaded}{}
\providecommand{\fnRows}{112}
\providecommand{\fnCode}{62}
\providecommand{\fnDisagree}{38}
\providecommand{\fnKept}{9}
\providecommand{\fnOtherClass}{29}
\providecommand{\fnParse}{13}
\providecommand{\fnAbsent}{8}
\providecommand{\fnSkew}{4}
\providecommand{\fnInternal}{4}
\providecommand{\fnPenSettled}{96}
\providecommand{\fnPenHandled}{11}
\providecommand{\fnPenClosed}{4}
\providecommand{\fnPenOther}{7}
\providecommand{\fnPenRollout}{0}
\providecommand{\fnTrials}{39}
\providecommand{\fnTrialsMedian}{15}
\providecommand{\fnTrialsMode}{10}
\providecommand{\fnTrialsModeN}{16}
\providecommand{\fnTotals}{24}
\providecommand{\fnTotalsMin}{12}
\providecommand{\fnTotalsMax}{750}
\providecommand{\fnUnclear}{7}
}{}}{}
\makeatother

\subsection{Reporting practice}
\label{subsec:reporting}

\begin{figure}[!t]
\centering
\IfFileExists{figs/fig_reporting.tex}{\resizebox{\columnwidth}{!}{% generated by tools/make_tikz_figures.py; do not edit
\begin{tikzpicture}[
  font=\footnotesize,
  bar/.style={draw=none,fill=black!62},
  barlight/.style={draw=none,fill=black!22},
  baracc/.style={draw=none,fill=accent},
  box/.style={draw=black!45,line width=0.3pt,inner sep=3pt,align=left},
  ghost/.style={draw=black!35,line width=0.3pt,dash pattern=on 1.4pt off 1.4pt,inner sep=3pt,align=left},
  lnk/.style={draw=black!38,line width=0.28pt},
  lbl/.style={font=\scriptsize},
  num/.style={font=\scriptsize\color{black!55}},
]

\node[lbl,anchor=east] at (0,-0.00) {success criterion stated};
\fill[barlight] (0.12,-0.10) rectangle (6.32,0.10);
\fill[bar] (0.12,-0.10) rectangle (5.54,0.10);
\node[num,anchor=west] at (6.44,-0.00) {98 (88\%)};
\node[lbl,anchor=east] at (0,-0.44) {real-robot experiment};
\fill[barlight] (0.12,-0.54) rectangle (6.32,-0.34);
\fill[bar] (0.12,-0.54) rectangle (5.05,-0.34);
\node[num,anchor=west] at (6.44,-0.44) {89 (79\%)};
\node[lbl,anchor=east] at (0,-0.88) {real trial count stated};
\fill[barlight] (0.12,-0.98) rectangle (6.32,-0.78);
\fill[bar] (0.12,-0.98) rectangle (4.00,-0.78);
\node[num,anchor=west] at (6.44,-0.88) {70 (62\%)};
\node[lbl,anchor=east] at (0,-1.32) {code released};
\fill[barlight] (0.12,-1.42) rectangle (6.32,-1.22);
\fill[bar] (0.12,-1.42) rectangle (3.55,-1.22);
\node[num,anchor=west] at (6.44,-1.32) {62 (55\%)};
\node[lbl,anchor=east] at (0,-1.76) {unseen-object count stated};
\fill[barlight] (0.12,-1.86) rectangle (6.32,-1.66);
\fill[baracc] (0.12,-1.86) rectangle (2.28,-1.66);
\node[num,anchor=west] at (6.44,-1.76) {39 (35\%)};
\node[lbl,anchor=east] at (0,-2.20) {contact or penetration handled};
\fill[barlight] (0.12,-2.30) rectangle (6.32,-2.10);
\fill[baracc] (0.12,-2.30) rectangle (0.73,-2.10);
\node[num,anchor=west] at (6.44,-2.20) {11 (10\%)};
\draw[lnk] (0.12,-2.48) -- (6.32,-2.48);
\end{tikzpicture}}}{%
  \framebox[0.95\columnwidth]{\rule{0pt}{2.2cm}\textit{figs/fig\_reporting.tex pending}}}
\caption{What gets reported, and what does not: six shares of the 112 corpus method rows, each bar
drawn against the denominator the text gives, with the two rarest reporting practices marked in
colour.}
\label{fig:reporting}
\end{figure}

Of the 112 method rows in the corpus, 89 report a real-robot experiment, 22 do not and one row is
unsettled, which is 80 percent of the 111 the note settled. Among those 89, 70 state how many real
trials produced the headline number, 79 percent of them. The 89 is the denominator that belongs to
this statistic: the 22 rows with no real robot cannot state a real trial count, and counting them
as silent turns a definitional impossibility into a reporting failure. Thirty-nine rows state a
count of unseen test objects, 35 percent. Ninety-eight state how a rollout is scored, 88 percent.
Sixty-two released code and 46 did not, with four rows unsettled, 57 percent of the 108 the note
settled. \figurename~\ref{fig:reporting} draws these six shares, each against the denominator that
belongs to it. Every bar is a lower bound, for the reason the next paragraph sets out. Two of the
six are marked, because they are the pair a reader needs in order to compare any two methods:
without a count of unseen objects a success rate has no generalisation denominator, and without
contact or penetration handling it says nothing about whether the hand passed through the object.
They are also the two the field reports least.

\textbf{These are counts of what this survey captured, not of what papers reported, and every one
is a floor.} A structured row holds a scalar. A paper that reports ten trials on each of nine
tasks, or a scoring rubric instead of a threshold, or a count spread over four tables, has nothing
the extraction can reduce to one integer, so it produces a null, and a null is then
indistinguishable from a paper that said nothing. The bias runs one way: every miss converts a
reporting paper into a silent one, and the survey's argument is that the field reports badly, so
the artefact flatters the argument. Section~\ref{sec:train:code} makes the same disclosure about
the paper/code count, subtracting the 13 disagreements that are limitations of this survey's own
parsing before declaring which number to quote, and the coverage statistics above need it more.

The nulls behind those six shares were therefore audited by hand against the notes they came from,
and every number above is post-audit. Appendix~\ref{app:method} gives what that audit recovered and
what it could not: the counts moved by tens of rows, 19 trial counts and 4 criteria remain
genuinely unsettled, and every bar in \figurename~\ref{fig:reporting} is still a lower bound.

The remaining gap is the one that matters. Nineteen of the 89 papers with a real robot never say
how many times they ran it. A percentage with no denominator cannot be given an interval, so it
cannot be compared with anything.

Where the denominator is stated it is small, and it is not one quantity. Some stored counts are
per-cell, meaning ten trials on each task, or twenty per condition, or five per object. Others are
grand totals over every cell. The rows now record which of the two each value is, and the two
distributions are quoted separately. The 39 per-cell counts run from 5 to 100 with a median of 15
and quartiles at 10 and 20. The modal cell is 10 trials, in 16 rows, then 20, in 12. The 24 grand
totals run from 12 to 750 with a median of 110. Pooling the two gives a median of 20 and a range
of 5 to 1287, and that pooled figure is the one an earlier draft of this section quoted. It
describes nothing, because Hora \cite{hora_2022}'s 240 and Visual Dexterity
\cite{visual_dexterity_2022}'s 20 are experiments of comparable size recorded on different bases.
Seven further counts could not be assigned a basis at all. \figurename~\ref{fig:trials} is the
per-cell distribution.

% ---------------------------------------------------------------------------
% The per-cell real-trial counts. The grand totals are a different quantity and
% are not pooled into it. Written by tools/make_finding_figures.py.
% ---------------------------------------------------------------------------
\begin{figure}[!t]
  \centering
  \IfFileExists{figs/fig_trials.tex}{\resizebox{\columnwidth}{!}{% generated by tools/make_finding_figures.py; do not edit
\begin{tikzpicture}[
  font=\footnotesize,
  bar/.style={draw=none,fill=black!62},
  barmid/.style={draw=none,fill=black!42},
  barlight/.style={draw=none,fill=black!20},
  baracc/.style={draw=none,fill=accent},
  seg/.style={draw=white,line width=0.5pt},
  lnk/.style={draw=black!38,line width=0.28pt},
  ghost/.style={draw=black!45,line width=0.3pt,dash pattern=on 1.2pt off 1.2pt},
  lbl/.style={font=\scriptsize},
  small/.style={font=\scriptsize\color{black!55}},
  num/.style={font=\scriptsize\color{black!55}},
]

\fill[bar] (0.00,0) rectangle (0.56,0.13);
\node[num,anchor=south] at (0.28,0.15) {1};
\draw[lnk] (0.56,0) -- (0.56,-0.06);
\node[lbl,anchor=north] at (0.56,-0.08) {5};
\fill[bar] (0.70,0) rectangle (1.26,2.10);
\node[num,anchor=south] at (0.98,2.12) {16};
\draw[lnk] (1.26,0) -- (1.26,-0.06);
\node[lbl,anchor=north] at (1.26,-0.08) {10};
\fill[bar] (1.40,0) rectangle (1.96,0.53);
\node[num,anchor=south] at (1.68,0.55) {4};
\draw[lnk] (1.96,0) -- (1.96,-0.06);
\node[lbl,anchor=north] at (1.96,-0.08) {15};
\fill[bar] (2.10,0) rectangle (2.66,1.58);
\node[num,anchor=south] at (2.38,1.60) {12};
\draw[lnk] (2.66,0) -- (2.66,-0.06);
\node[lbl,anchor=north] at (2.66,-0.08) {20};
\fill[bar] (2.80,0) rectangle (3.36,0.26);
\node[num,anchor=south] at (3.08,0.28) {2};
\draw[lnk] (3.36,0) -- (3.36,-0.06);
\node[lbl,anchor=north] at (3.36,-0.08) {25};
\draw[lnk] (4.06,0) -- (4.06,-0.06);
\node[lbl,anchor=north] at (4.06,-0.08) {30};
\fill[bar] (4.20,0) rectangle (4.76,0.13);
\node[num,anchor=south] at (4.48,0.15) {1};
\draw[lnk] (4.76,0) -- (4.76,-0.06);
\node[lbl,anchor=north] at (4.76,-0.08) {35};
\fill[bar] (4.90,0) rectangle (5.46,0.13);
\node[num,anchor=south] at (5.18,0.15) {1};
\draw[lnk] (5.46,0) -- (5.46,-0.06);
\node[lbl,anchor=north] at (5.46,-0.08) {40};
\draw[lnk] (6.16,0) -- (6.16,-0.06);
\node[lbl,anchor=north] at (6.16,-0.08) {45};
\fill[bar] (6.30,0) rectangle (6.86,0.13);
\node[num,anchor=south] at (6.58,0.15) {1};
\draw[lnk] (6.86,0) -- (6.86,-0.06);
\node[lbl,anchor=north] at (6.86,-0.08) {50};
\fill[bar] (7.34,0) rectangle (7.90,0.13);
\node[num,anchor=south] at (7.62,0.15) {1};
\draw[lnk] (7.62,0) -- (7.62,-0.06);
\node[lbl,anchor=north] at (7.62,-0.08) {$>$50};
\draw[lnk] (-0.06,0) -- (7.96,0);
\node[lbl,anchor=north] at (3.95,-0.42) {real trials per cell, in bins of five};
\draw[draw=accent,line width=0.4pt,dash pattern=on 1.4pt off 1.2pt] (1.96,-0.02) -- (1.96,2.44);
\node[anchor=west,font=\scriptsize\color{accent}] at (2.04,2.34) {median 15};
\node[small,anchor=north west,align=left] at (-0.06,-0.66) {39 rows state a per-cell count; the mode is 10, in 16 of them\\24 state a grand total instead, 12 to 750, and are not pooled into this\\7 more state a count on a basis that could not be read};
\end{tikzpicture}
}}{%
    \framebox[0.95\columnwidth]{\rule{0pt}{2.2cm}\textit{figs/fig\_trials.tex pending}}}
  \caption{Where a denominator is stated it is small: the \fnTrials{} per-cell real-trial counts,
  binned by fives, median \fnTrialsMedian{} marked. The \fnTotals{} grand totals are a separate
  quantity, counted beneath the axis rather than pooled in.}
  \label{fig:trials}
\end{figure}

The per-cell count is the one a Wilson interval attaches to. For $s$ successes in $n$ trials, with
$\hat{p} = s/n$ and $z = z_{1-\alpha/2} = 1.96$ at the 95 percent level, the bounds of that
interval are
\begin{equation}
\frac{2n\hat{p} + z^{2} \;\pm\; z\sqrt{4n\hat{p}\,(1-\hat{p}) + z^{2}}}{2\,(n + z^{2})},
\label{eq:wilson}
\end{equation}
which, unlike the normal approximation, stays inside $[0,1]$ and does not collapse to zero width at
$\hat{p} = 0$ or $\hat{p} = 1$. At 20 trials a reported 60 percent carries a 95 percent Wilson
interval of 39 to 78 percent, and a reported 80 percent an interval of 58 to 92 percent.
Overlapping intervals are not a test, so the point is made with one: 12 of 20 against 16 of 20 is
$z = 1.38$, $p = 0.17$, and the difference is not established. Two methods separated by 20 points
at 20 trials, the larger of the two modal cell sizes, are not separated at all.

The share of papers stating a count has risen, from 39 of the 65 rows before 2025, 60 percent, to
31 of the 47 rows from 2025 and 2026, 66 percent. The counts themselves have not. The median
stated count is 20 in 2024, 22.5 in 2025 and 20 in 2026. Before 2024 the per-year medians rest on
one to seven observations and should not be read as a trend.

The denominators also sit on different hardware. The 103 method rows that name their own hand give
78 distinct hand strings between them, and this survey applies no normalisation to those strings,
so 78 is a count of strings and not of hand designs. Matching on the string, Allegro appears in
35, Shadow in 21, Inspire in 19 and LEAP in 12. A success rate on a 16-degree-of-freedom Allegro
and a success rate on a 6-actuator Inspire hand are not measurements of the same thing.

Nor are the criteria. DexVerse
\cite{dexverse_2026} counts PickCube a success when the cube is ``lifted at least 0.20 m above its
resetting height''. Bench2Dex \cite{bench2dex_2026} requires its terminal predicate to hold for a
continuous dwell time of 0.5 s, to reject transient contacts. The Colosseum \cite{colosseum_2024}
counts an episode successful ``if the model completes the task fully''. DexTrack
\cite{dextrack_2025} reports every success rate as a pair under two threshold sets, which on GRAB
gives 46.70 and 65.48 percent for the same rollouts. Of the 14 rows that still state no criterion,
ten have no success predicate at all. They report radians rotated or time-to-fall and never define
a success, which is a fact about the paper rather than a gap in this survey. And four are
unsettled by the note. Of the 19 criteria the audit recovered, all but one are a rubric, a staged
partial credit, a judgement by eye or a deferral to a benchmark's own definition rather than a
threshold, and exactly one, \cite{pistar06_2025}, states a verbatim numeric one;
Appendix~\ref{app:method} counts the four kinds. A rubric is a milder failure than silence and a worse one than a
threshold, because it is reproducible inside a lab and not across two.

\textbf{The axes that matter.}\label{subsec:axes}
Seven quantities dissociate in the published data, so they have to be reported separately.

\textbf{Task success.} Binary success discards the difference between near-misses and inaction.
Snyder et al. \cite{beyond_binary_success_2026} put it plainly: ``a policy that completes 90\% of
the task is clearly better than a policy that is frozen the whole time, yet their success rates
would be identically 0\%''. In \cite{kress_gazit_policy_eval_2024} policy C scores 17 percent
overall on the pancake task while picking up the spatula and flipping the pancake in 23 of 23
attempts.

\textbf{Robustness to perturbation.} The Colosseum \cite{colosseum_2024} measures a 30 to 50
percent success drop under single perturbation factors and at least 75 percent under all 14
together. In \cite{bench2dex_2026} GR00T N1.5 leads the matched condition at 48.5 percent and
falls to 19.8 percent under combined shift, while $\pi_{0.5}$ goes from 27.3 to 19.7 and retains
the most at 72.1 percent. The ranking at the anchor is not the ranking under shift.

\textbf{Generalisation to unseen objects.} Only 39 rows state a count and the median is 11
objects.

\textbf{Physical plausibility of the contact.} Eleven of the 96 rows whose contact handling the
note settled address it, 11 percent, with 16 rows unknown. Section~\ref{subsec:plausibility} takes
them apart.

\textbf{Sample and wall-clock cost.} Thirty-one of 112 rows state a parallel environment count and
18 a simulated episode count. RoboPianist \cite{robopianist_2023} is the exception, at 5 million
samples per song and roughly 5 hours per run on four Tesla K80 GPUs.

\textbf{Real-robot transfer.} Simulated rank order is not real rank order. AutoEval
\cite{autoeval_2025} scores Open-$\pi_0$ on put-eggplant-in-sink at 6 of 50 in SIMPLER and 47 of
50 on the real WidowX. Badithela et al. \cite{suresim_2025} state the limit directly: ``the
simulation-to-real gap precludes rigorous statistical inferences about real-world outcomes from
simulation results alone''.

\textbf{Reproducibility.} 62 rows released code that could be parsed against the paper, and 38 of
the 112 rows record a disagreement of some kind between the paper and that code. All 38 released
code, so the raw rate among code-releasing rows is 61 percent. That raw rate is not the finding,
because the 38 are not one thing. Section~\ref{sec:train:code} classifies them: 9 contradictions, 13 limitations
of this survey's own parsing, 8 components never released, 4 version skews and 4 inconsistencies
internal to a paper. Only the contradictions are a finding about the work rather than about this
survey, so 9 of 62 code-releasing rows, which is 15 percent, is the figure this section and
Table~\ref{tab:protocol} use. PhysHOI \cite{physhoi_2023} is the clearest of the 9. It lists a
non-zero object-orientation weight for GRAB in its Table 4, and the reward function it released
hard-sets that orientation error to zero, so the reward that produced the published numbers
tracked the object in position only.

\subsection{Physical plausibility}
\label{subsec:plausibility}

% ---------------------------------------------------------------------------
% The second finding as a picture: the funnel that ends at zero. Written by
% tools/make_finding_figures.py from corpus/rows.
% ---------------------------------------------------------------------------
\begin{figure}[!t]
  \centering
  \IfFileExists{figs/fig_penetration.tex}{%
    \resizebox{\columnwidth}{!}{% generated by tools/make_finding_figures.py; do not edit
\begin{tikzpicture}[
  font=\footnotesize,
  bar/.style={draw=none,fill=black!62},
  barmid/.style={draw=none,fill=black!42},
  barlight/.style={draw=none,fill=black!20},
  baracc/.style={draw=none,fill=accent},
  seg/.style={draw=white,line width=0.5pt},
  lnk/.style={draw=black!38,line width=0.28pt},
  ghost/.style={draw=black!45,line width=0.3pt,dash pattern=on 1.2pt off 1.2pt},
  lbl/.style={font=\scriptsize},
  small/.style={font=\scriptsize\color{black!55}},
  num/.style={font=\scriptsize\color{black!55}},
]

\node[lbl,anchor=south west] at (0,0.04) {\textbf{112}~~method rows in the corpus};
\fill[black!8] (0,0.00) rectangle (7.60,-0.26);
\fill[barlight] (0,0.00) rectangle (7.60,-0.26);
\node[lbl,anchor=south west] at (0,-0.56) {\textbf{96}~~of them whose note settles the question};
\fill[black!8] (0,-0.60) rectangle (7.60,-0.86);
\fill[barmid] (0,-0.60) rectangle (6.51,-0.86);
\node[lbl,anchor=south west] at (0,-1.16) {\textbf{11}~~of those 96 address interpenetration in any form};
\fill[black!8] (0,-1.20) rectangle (7.60,-1.46);
\fill[bar] (0,-1.20) rectangle (0.75,-1.46);
\node[lbl,anchor=south west] at (0,-1.76) {\textbf{4}~~of those 11 inside a closed-loop policy};
\fill[black!8] (0,-1.80) rectangle (7.60,-2.06);
\fill[baracc] (0,-1.80) rectangle (0.27,-2.06);
\fill[wash] (0,-2.52) rectangle (7.60,-3.46);
\draw[line width=0.3pt,draw=black!30] (0,-2.52) rectangle (7.60,-3.46);
\node[anchor=west,font=\fontsize{22}{22}\selectfont\bfseries\color{accent}] at (0.18,-2.99) {0};
\node[anchor=west,align=left,font=\scriptsize] at (1.00,-2.99) {report a penetration number for the rollouts\\of their own trained policy};
\end{tikzpicture}
}}{%
    \framebox[0.95\columnwidth]{\rule{0pt}{2.2cm}\textit{figs/fig\_penetration.tex pending}}}
  \caption{Interpenetration handling among all \fnRows{} method rows, the settled \fnPenSettled{}
  shown against the light track behind them, split by closed-loop policy versus reference-only
  scoring.}
  \label{fig:penetration}
\end{figure}

The quantity most specific to a hand is the one closed-loop policies do not record, and
\figurename~\ref{fig:penetration} is the funnel that narrows to zero. Eleven method rows handle
interpenetration in any form: three penalise it, three measure it, five constrain it, 11 percent
of the settled rows. The denominator is 96, not 112, because the
contact-handling field is null for 16 rows, and a null there means the note did not settle the
question, not that the paper ignored penetration.

That eleven is not a claim that penetration goes unmeasured in general, and reading it that way
would be wrong. Outside closed-loop control the quantity is a standard comparative column, and has
been one for years in grasp synthesis and in hand-object reconstruction. Four rows of this corpus
show the practice. BiDexGrasp \cite{bidexgrasp_2026} prints a penetration depth beside a prior
method's, BimanGrasp \cite{bimangrasp_2024} fails any grasp that exceeds 1.5 mm of total
penetration, TopoRetarget \cite{toporetarget_2026} reports a maximum penetration and a share of
frames past 2 mm against a baseline retargeter, and OakInk \cite{oakink_2022} scores a dataset
split on penetration depth, solid intersection volume and simulation displacement. The finding is
narrower than the field and concerns learned closed-loop control: all eleven score a pose or a
reference trajectory, four of them are closed-loop policies, and we found none that reports the
measurement for rollouts of its own trained policy.

Where in the pipeline those eleven act is the reference-versus-rollout split that
Section~\ref{sec:intro} takes from \cite{zhao_dexhand_survey_2026}. A reference is a pose or a
trajectory scored before execution, and a rollout is what the trained policy actually did. What
this section supplies on that axis is a measurement method and a count, and it does not supply a
threshold. The count is the eleven of 96 above, with four closed-loop policies inside it and none
we found reporting a number for its own rollouts, which is \figurename~\ref{fig:penetration}. The
method is the plausibility row of Table~\ref{tab:protocol}: maximum and mean penetration depth
over the evaluation rollouts, on a dense surface sample, computed by code that never entered the
reward or the termination rule. The threshold is borrowed, and
Section~\ref{subsec:what_would_be_true} says from where and why it does not bind. Six of the
eleven are grasp synthesisers or trajectory optimisers, namely
\cite{bidexgrasp_2026,bimangrasp_2024,deximit_2026,pang_global_planning_2022,toporetarget_2026}
and \cite{unidexgrasp_2023}, and Castro et al. \cite{castro_sap_contact_2021} give a contact model
rather than a controller. That leaves four closed-loop policies in the whole corpus:
\cite{clutterdexgrasp_2025,dexmachina_2025,dextrack_2025} and \cite{teledexter_2026}.

The reason the number is four is a measurement trap. A quantity a policy optimises cannot also
judge it, because the policy learns the measure rather than the property the measure stands for.
PhysHOI \cite{physhoi_2023} documents the failure in the clean direction. Its kinematic imitation
reward was maximised by not touching the object at all, because contact perturbed the reference
trajectory: ``the policy may learn not to touch the object and falls into a local optimal''. A
contact-graph reward closed the hole, and that reward reads a force threshold rather than
geometry, so it constrains contact presence and says nothing about penetration depth.

DexTrack \cite{dextrack_2025} shows the trap in the other direction. It defines a maximum
hand-object penetration depth over all frames in its Appendix B, and applies it only to the input
kinematic references, as one component of a reference-quality score. No penetration number appears
for any of its own rollouts, in simulation or on the LEAP hand. Tolerance of the failure is then
reported as a result: ``Despite severe hand-object penetrations in Figure 4c and Figure 4a, the
hand still interacts effectively with the object, highlighting the resilience of our tracking
controller''.

TopoRetarget \cite{toporetarget_2026} is the strongest case in the corpus and still stops one step
short on the same reference-versus-rollout line. It constrains penetration during retargeting with
a 1 mm soft tolerance and a 30 mm hard bound, and it reports two numbers on 25 ContactPose grasps:
a maximum penetration of 1.07 mm and 0.00 percent of frames above 2 mm, against 22.22 mm and 96
percent of frames for its GeoRT baseline. Then a PPO controller tracks those references, and its
four reward terms and its five termination criteria govern object pose, link position, joint error
and action smoothness, never penetration. The constrained quantity is the reference, and the
rollout is not re-measured.

Definitions are not shared either. GRAB \cite{grab_2020} estimates contact by proximity, because
``contact cannot be directly observed'', with a 4.5 mm tolerance, and reports that ``\,`Use'
grasps have $3.25 \pm 0.68$ mm average penetration'', without saying whether 3.25 mm is a maximum,
a mean or a median. OakInk \cite{oakink_2022} supplies the fullest published vocabulary:
penetration depth, solid intersection volume and simulation displacement. Its Table 3 scores the
GRAB GrabNet split at 2.53 cm penetration depth, against GRAB's own 3.25 mm. The two differ by a
factor of about eight, and neither source states its distance function precisely enough to
reconcile them, and the two are not scored over the same grasps either, since GRAB's figure is
over its own captured ``use'' grasps and OakInk's is a model's output on the GrabNet split, so a
difference of population and a difference of definition are confounded in the same ratio.

The analytic tradition scored a grasp without simulating it, and Ferrari and Canny's measure
\cite{ferrari_canny_1992} and Roa and Su\'arez's review \cite{roa_suarez_grasp_quality_2015} are
its wrench-space reference points. Neither could be obtained. The first DOI fetch returned HTTP
202 and the second a Springer JavaScript interstitial, so both are cited by metadata only and no
definition here rests on their contents. The learned literature has not replaced that tradition
with anything it measures on its own rollouts. Penetration is a quantity between meshes, so it
needs a simulator or a mesh reconstruction, and the protocol below treats it as a simulation-only
axis.

\textbf{What would close it.} The gap is specific to learned closed-loop control rather than general,
and it is a choice rather than a capability. What would close it is a maximum and a mean penetration
depth over the evaluation rollouts, on a dense surface sample, computed by a measure the policy never
optimised.

\subsection{Statistical practice, and a proposed protocol}
\label{subsec:statistics}

Kress-Gazit et al. \cite{kress_gazit_policy_eval_2024} give the field's reference protocol, and it
prescribes process, not numbers. Write the success criteria before the run and have someone other
than their author score the runs. Match initial conditions with image overlays. Interleave the
policies blind within one session. Report counts rather than percentages, alongside the initial
conditions and the failure modes. Use a posterior over the Bernoulli parameter instead of a point
estimate. It prescribes no minimum trial count anywhere and no frequentist confidence-interval
width anywhere. Its own example report uses 10 initial conditions with two runs each, 20
evaluations per policy, and its worked case shows what 20 buys. Pancake success of 15 of 18
against 11 of 17, nominally 83 against 65 percent, leaves a 0.11 posterior probability that the
worse policy is actually better. At 150 of 180 against 110 of 170 the same rates separate.

The TRI LBM Team \cite{lbm_careful_examination_2025} supplies the missing numbers by fiat rather
than derivation: 50 rollouts per task per policy per condition on hardware, 200 in simulation,
blind, with randomised policy order inside per-initial-condition bundles. It replaces confidence
intervals with Beta-posterior violins and corrects its pairwise tests under Bonferroni, ``unless
otherwise noted''. Its own warning is the strongest sentence in this literature: ``there is
significant risk that many robotics papers are measuring statistical noise due to insufficient
statistical power''.

The rest fall short of their own advice. RoboArena \cite{roboarena_2025} runs 612 double-blind
pairwise comparisons across seven institutions and 4284 rollouts, and reports no confidence
intervals and no per-policy trial counts. The Colosseum \cite{colosseum_2024} evaluates 235 test
sets at 25 episodes each with ``one training seed and one evaluation seed'' per baseline and gives
no intervals in simulation. AutoEval \cite{autoeval_2025} runs 50 trials per policy per task and
calls $\pm 10$ percent ``the natural variance of robot evaluations'', without giving the formula
behind the intervals it plots.

Two papers do give usable numbers. Badithela et al. \cite{suresim_2025} pair real and simulated
trials and de-biases the simulation with a rectifier, saving 20 to 25 percent of hardware trials
at paired correlations of roughly 0.6 to 0.7, and nothing at all at a correlation near zero. The
25 percent figure is the better of its two reported settings, the DP case at $\rho = 0.702$. Its
decision rule is exact: combining helps only when the rectifier variance is below the variance of
the real evaluations. Snyder et al. \cite{beyond_binary_success_2026} give the largest saving. On
the LBM rubrics their sequential test on graded scores cuts simulated evaluation by about 70
percent and hardware by about 45 percent, 286 trials against a nominal 500, with per-task
decisions landing in 12 to 36 paired trials. On RoboArena's data a 30-point gap on continuous
progress scores reaches significance in 18 trials, while a 20-point gap on binary success needs
about 80.

\textbf{A proposed protocol.}\label{subsec:protocol}
Every count below is printed by \path{tools/make_eval_tables.py} in its derivation mode, and each
axis is derived for the statistic that axis actually reports: a single rate takes a Wilson
half-width, a matched comparison takes McNemar, a ratio takes the standard error of the log ratio,
a correlation takes the Fisher-z interval.

Appendix~\ref{app:protocol} works the derivation; what it concludes is the following, and every
count in Table~\ref{tab:protocol} is one of these. A single reported rate takes 100 real rollouts
per cell, because 93 is where the worst-case Wilson half-width of \eqref{eq:wilson} falls to 10
points and 10 points is the coarsest width at which a rate is worth printing. A comparison is a
different question and a harder one. The design is paired, so the count follows McNemar and depends
on the discordance rate rather than on the two rates alone, and separating 50 from 70 percent at
80 percent power takes 57 matched pairs per arm at the assumed discordance of 0.3, which is 114
rollouts against the 186 two independent arms would need. Simulated cells take 200 episodes for a
6.9-point half-width. Perturbation axes are screened at 40 per axis, which ranks the axes and
cannot establish that any one of them hurt, and certified at 101 per arm. Unseen objects resample
the object rather than the trial, at 20 objects and 5 trials each, and a 10-point claim about an
object distribution needs about 93 objects. The transfer axis takes the 100 real trials paired to
100 simulated, and about 457 pairs to separate two correlations 0.1 apart. One hundred is a cap and
not a bill, because a sequential test on a graded score reached its decision in 12 to 36 paired
trials in \cite{beyond_binary_success_2026}, and a cell that stops early reports a confidence
sequence rather than a Wilson interval, never both. The appendix also keeps the counts an earlier
draft of this section quoted and this one withdrew, since the test they were computed for was the
wrong one.

\IfFileExists{tables/table8_protocol.tex}{% alias for table7_protocol.tex; emits once however many sections input it
\ifcsname dxdone@table7_protocol\endcsname\else
  \expandafter\gdef\csname dxdone@table7_protocol\endcsname{}%
  % generated by tools/make_tex_tables.py -- do not edit
\begin{table*}[!t]
\centering
\caption{The proposed evaluation protocol, and this survey's own contribution. No cell is empty: every cell is a proposal.}
\label{tab:protocol}
\footnotesize
\setlength{\tabcolsep}{3pt}
\setlength{\emergencystretch}{1.5em}
\begin{tabular}{>{\raggedright\arraybackslash\hspace{0pt}}p{62pt}>{\raggedright\arraybackslash\hspace{0pt}}p{128pt}>{\raggedright\arraybackslash\hspace{0pt}}p{108pt}>{\raggedright\arraybackslash\hspace{0pt}}p{170pt}}
\toprule
\rowcolor{wash} \hdr{axis} & \hdr{what is measured} & \hdr{minimum trials} & \hdr{reported with it} \\
\midrule
Task success & fraction of episodes meeting a criterion written before the run, and a graded rubric score in $[0,1]$ & 100 real and 200 sim per policy, task and condition for an absolute rate at a 10-point interval; a comparison alone needs 57 matched pairs per arm & raw counts, the criterion verbatim, the rubric, the initial-condition protocol, and an interval whose kind is named \\
Robustness & success under each perturbation axis, and the ratio to the unperturbed anchor & 40 per axis per policy to rank the axes; 101 per axis and 101 at the anchor before any claim that a named axis hurt & both absolute rates with their Wilson intervals, the axis ranges, and the ratio without an interval unless the certifying count was run \\
Unseen objects & mean over held-out objects of per-object success & 20 objects at 5 trials to screen; 93 objects for a 10-point claim & the object list, the per-object rates, a bootstrap interval over objects, and the within-object trial count \\
Physical plausibility & maximum and mean hand-object penetration depth per rollout, and the fraction of frames above 2\,mm & 100 rollouts, simulation only & an interval over rollouts, the sample density, the threshold, the solver's depenetration settings, and at least one rendered rollout \\
Sample and wall-clock cost & environment steps and wall-clock to the checkpoint that produced the headline number & 3 training seeds & the range over seeds, and the statement that 3 seeds is a range and not an interval \\
Real-robot transfer & paired real and simulated outcomes on matched initial conditions, their correlation, and the rectifier variance against the real variance & the 100 real trials paired to 100 sim; about 457 pairs to separate two correlations 0.1 apart & the correlation with its Fisher-z interval, the variance ratio with its bootstrap interval, and which decision the count supports \\
Reproducibility & the released config and checkpoints against the paper's stated objective & not a trial count & the file and line of every disagreement, or an explicit statement that none was found \\
\bottomrule
\end{tabular}
\end{table*}
%
\fi
}{%
  \begin{table}[!t]\centering
    \caption{Pending: the protocol table, generated by the table generator.}
  \label{tab:protocol}
  \framebox[0.95\columnwidth]{\rule{0pt}{1.4cm}}
  \end{table}}

\subsection{The results matrix, and what filling it would cost}
\label{subsec:matrix}

The rows are the 12 most-mentioned dexterous-hand policy methods in the corpus, and the rule is
the one the table's generator implements, stated here in the same words. A candidate is a method
row that names a hand; the hand string must not contain ``parallel'' or ``gripper''; it must carry
at least one paradigm tag that produces a closed-loop policy and must not be a teleoperation
system, because an interface is scored on latency and operator effort rather than on a policy's
success rate. And its name must be at least four characters, so that a short string does not match
everything. Candidates are then scored by the number of other corpus papers whose parsed text
contains the name, and the top 12 by count, ties broken by key, are the rows.

Two corrections changed that ranking. The match is on a whole word. Under the bare substring test
an earlier version used, ``UniDex'' matched inside ``UniDexGrasp'' and ``UniDexGrasp++'', and
\cite{unidex_2026}. A 2026 paper. Sat sixth in a ranking over a corpus written mostly before it,
on 34 mentions that belonged to a different work. As a whole word it has 3 and it is not in the
table. And the interface rule is now applied to every row that carries the tag rather than only to
rows that carry nothing else, which drops \cite{anyteleop_2023} at 34 mentions, \cite{dime_2022}
at 28 and \cite{holo_dex_2022} at 23, along with \cite{dexpilot_2020}, which the earlier prose
already excluded by hand.
%% The counts themselves, generated with the table by tools/make_tex_tables.py. They used to be set
%% under the float in six-point type; they are one sentence of the paragraph that states the rule
%% they follow from.
% generated by tools/make_tex_tables.py -- do not edit
% The mention counts behind the ranking of tables/table8_matrix.tex, as a sentence for the
% body of Section VII, which is where the ranking rule they follow from is stated.
The counts, on whole-word matches over \path{papers/md}, are OpenAI~64; DexMV~50; DAPG~49; DeXtreme~44; DexCap~40; Hora~34; Visual Dexterity~32; Yin et al.~27; UniDexGrasp~27; PDDM~25; HATO~23; DexPoint~22.


The ranking is not one quantity even so. A method's name is taken from the first line of its note,
which yields an acronym for some works and a full title for others, and a title is matched mostly
inside reference lists while an acronym is matched in running text. Those have different base
rates, so the table marks which kind each row was matched on and the two kinds are not comparable
with each other. Mention counts are counts of mentions and not of use, as
Appendix~\ref{app:method} records.

Every cell is empty. This survey re-ran nothing, and no cell can be filled at the denominator
Table~\ref{tab:protocol} asks for. DeXtreme \cite{dextreme_2022} comes closest and is the reason
the claim is stated that narrowly: it reports a criterion, a trial count and an interval. Object
orientation within 0.4 rad of target, $27.8 \pm 19.0$ average consecutive successes with the $\pm$
a 90 percent confidence interval. On 10 trials. The method tabulation is not a counter-example
either, though it looks like one: it carries trial-count, unseen-object, penetration and
code-release columns for all 112 method rows, including all 12 of these. It records what each
method reported. Table~\ref{tab:matrix} asks instead for what Table~\ref{tab:protocol} defines. A
value with an interval, a stated denominator and a criterion written before the run. And no
reported value meets that. The first row of Table~\ref{tab:matrix} is a worked example so that the
format of a cell is unambiguous. Every number in it is fabricated and labelled as such.

\IfFileExists{tables/table9_matrix.tex}{% alias for table8_matrix.tex; emits once however many sections input it
\ifcsname dxdone@table8_matrix\endcsname\else
  \expandafter\gdef\csname dxdone@table8_matrix\endcsname{}%
  % generated by tools/make_tex_tables.py -- do not edit
\begin{table*}[!t]
\centering
\caption{The matrix, for someone else to fill. 12 rows and 84 cells, all 84 of them empty. An empty cell here is not \na: it is a measurement nobody has made, which is the point of the table, and no cell in it has a source to be missing. The first row is a worked example and every number in it is fabricated. Rows are the 12 most-mentioned dexterous-hand policy methods in the corpus, by the rule of Section~VII, scored on whole-word matches over \texttt{papers/md}: \texttt{openai\_\allowbreak dexterity\_\allowbreak 2018}~64; \texttt{dexmv\_\allowbreak 2021}~50; \texttt{dapg\_\allowbreak 2017}~49; \texttt{dextreme\_\allowbreak 2022}~44; \texttt{dexcap\_\allowbreak 2024}~40; \texttt{hora\_\allowbreak 2022}~34; \texttt{visual\_\allowbreak dexterity\_\allowbreak 2022}~32; \texttt{rotating\_\allowbreak without\_\allowbreak seeing\_\allowbreak 2023}~27; \texttt{unidexgrasp\_\allowbreak 2023}~27; \texttt{pddm\_\allowbreak 2019}~25; \texttt{hato\_\allowbreak visuotactile\_\allowbreak 2024}~23; \texttt{dexpoint\_\allowbreak 2022}~22. A mention count is a count of mentions, not of use. $\ast$ marks a work matched on its title rather than on a short name; a title is matched mostly inside reference lists and a short name in running text, so the two kinds of count are not comparable. Each column is an axis of Table~\ref{tab:protocol} and is filled in the units that table asks for.}
\label{tab:matrix}
\footnotesize
\setlength{\tabcolsep}{3pt}
\setlength{\emergencystretch}{1.5em}
\begin{tabular}{>{\raggedright\arraybackslash\hspace{0pt}}p{66pt}>{\raggedright\arraybackslash\hspace{0pt}}p{55pt}>{\raggedright\arraybackslash\hspace{0pt}}p{55pt}>{\raggedright\arraybackslash\hspace{0pt}}p{55pt}>{\raggedright\arraybackslash\hspace{0pt}}p{55pt}>{\raggedright\arraybackslash\hspace{0pt}}p{55pt}>{\raggedright\arraybackslash\hspace{0pt}}p{55pt}>{\raggedright\arraybackslash\hspace{0pt}}p{55pt}}
\toprule
\rowcolor{wash} \hdr{method} & \hdr{task success} & \hdr{robustness} & \hdr{unseen objects} & \hdr{plausibility} & \hdr{cost} & \hdr{transfer} & \hdr{reproducibility} \\
\midrule
\emph{worked example} & 0.72 $\pm$ 0.09 (100) & lighting 0.41 (40), ratio 0.57 & 0.55 [0.41, 0.68], 20 obj & 3.1 / 0.8 mm, 0.12 & 1.2e9 steps, 46 GPU-h & 0.61 [0.47, 0.72] & \texttt{train.\allowbreak yaml:88} \\
\texttt{openai\_\allowbreak dexterity\_\allowbreak 2018}$\ast$ &  &  &  &  &  &  &  \\
\texttt{dexmv\_\allowbreak 2021} &  &  &  &  &  &  &  \\
\texttt{dapg\_\allowbreak 2017}$\ast$ &  &  &  &  &  &  &  \\
\texttt{dextreme\_\allowbreak 2022} &  &  &  &  &  &  &  \\
\texttt{dexcap\_\allowbreak 2024} &  &  &  &  &  &  &  \\
\texttt{hora\_\allowbreak 2022}$\ast$ &  &  &  &  &  &  &  \\
\texttt{visual\_\allowbreak dexterity\_\allowbreak 2022}$\ast$ &  &  &  &  &  &  &  \\
\texttt{rotating\_\allowbreak without\_\allowbreak seeing\_\allowbreak 2023}$\ast$ &  &  &  &  &  &  &  \\
\texttt{unidexgrasp\_\allowbreak 2023} &  &  &  &  &  &  &  \\
\texttt{pddm\_\allowbreak 2019}$\ast$ &  &  &  &  &  &  &  \\
\texttt{hato\_\allowbreak visuotactile\_\allowbreak 2024}$\ast$ &  &  &  &  &  &  &  \\
\texttt{dexpoint\_\allowbreak 2022} &  &  &  &  &  &  &  \\
\bottomrule
\end{tabular}
\end{table*}
%
\fi
}{%
  \begin{table}[!t]\centering
    \caption{Pending: the results matrix, generated by the table generator.}
  \label{tab:matrix}
  \framebox[0.95\columnwidth]{\rule{0pt}{1.4cm}}
  \end{table}}

\textbf{What would have to be true.}\label{subsec:what_would_be_true}
The bill comes first. Per policy and per task the protocol asks for 57 matched trials on the
anchor set, which is the paired count from \S\ref{subsec:protocol}, and 100 on the unseen-object
set, at 20 objects and 5 trials each, with robustness and plausibility absorbed by simulation. A
two-policy comparison on three tasks is then 342 real rollouts on the matched set and 600 on the
unseen-object set, 942 in all, and at one minute per rollout including the reset that is about 16
hours of robot time, before failed resets, repairs and scoring. The same bill computed from the
independent-arm count, which an earlier draft used, was 1200 rollouts and 20 hours. The pairing
removes about a fifth of it rather than the four fifths the pair counts suggest, because a pair is
two rollouts and only the anchor set is paired. AutoEval \cite{autoeval_2025} ran about 850
episodes in 24 hours on a WidowX with three human interventions, and had to pause 20 minutes every
6 hours once the motors overheated. A tendon-driven multi-finger hand is more fragile, so 16 hours
of rollouts is most of a week of calendar time.

That bill is large but not unprecedented. AutoEval \cite{autoeval_2025} records that evaluating
OpenVLA against its baselines took more than 2500 rollouts and more than 100 hours of human labour
across three institutions, and the TRI LBM Team \cite{lbm_careful_examination_2025} analysed about
1800 real rollouts across nine hardware stations. What is unprecedented is paying it for a single
dexterous-hand paper, where the modal per-cell count is 10 trials and the median per-cell count
15, and where the median stated count has been 20 or 22.5 in each of the last three years.

Four things would have to change. Reviewers would have to reward 57 matched trials on one task
over 20 unmatched trials on five, and nothing in the corpus suggests that is happening. The loop
would have to be automated, and AutoEval \cite{autoeval_2025} shows it is buildable for a gripper
at one to three hours of setup, while stating that it supports binary success only and no
robustness axes. Somebody would have to run the penetration measure on their own rollouts, which
is a choice rather than a capability: Section~\ref{sec:sim:engines} shows the depth is computable
from the poses and the meshes in a few lines of Warp, and IsaacGymEnvs already ships a task that
does it every step. An earlier draft of this section made the engine the barrier, and that claim
is withdrawn, because the released code refutes it. And the comparison would have to be
sequential, because the savings in \cite{beyond_binary_success_2026} are the only reason 100 is a
cap rather than a cost.

Four limits apply to the proposal itself. This survey re-ran no method, so every count in
Table~\ref{tab:protocol} is derived from an interval width, a power calculation or another paper's
measurement, and Table~\ref{tab:matrix} is empty because we filled no cell. The counts are
worst-case at $\hat{p} = 0.5$, so a method near 90 percent needs fewer trials for the same width
and a method near 50 percent needs all 100, and the paired count additionally rests on an assumed
discordance of 0.3, which no paper in this corpus reports. The perturbation axes are borrowed from
The Colosseum \cite{colosseum_2024} and SIMPLER \cite{simpler_2024}, which run parallel-jaw
grippers on rigid objects, where SIMPLER found physical parameters moved success rates by at most
15 percent. That is the sensitivity expected to grow with multi-finger contact, and nobody has
measured it. The 2 mm penetration threshold is taken from \cite{toporetarget_2026} with no
independent justification, and the captured human grasps in \cite{grab_2020} sit above it at 3.25
mm, which makes 2 mm a simulator convention rather than a physical bound. So the three things this
survey adds to the reference-versus-rollout axis are a count, a measurement method and a protocol
slot for them, and a threshold is not among them. It is borrowed from one paper and reported as
borrowed, and it will stay a convention until somebody measures penetration on rollouts across
hands and engines and finds a value that separates behaviour a physicist would accept from
behaviour they would not.

\textbf{The methodology literature has no dexterous hand in it.} Seven corpus rows are evaluation
protocols, a small denominator, and not one of them uses a dexterous hand: Badithela et al.
\cite{suresim_2025} run a parallel-jaw gripper and the other six state no hand. Everything proposed
above is therefore assembled from work on grippers and on whole-arm tasks. What would close that, and
close it cheaply, is to run the Table~\ref{tab:protocol} protocol once, on one in-hand reorientation
task with one 16-DoF hand, and release the rollouts as the first row of Table~\ref{tab:matrix}.